\documentclass{article}
\usepackage{amsmath}
\usepackage{amssymb}
\usepackage{amsthm}
\usepackage{mathtools}
\usepackage{geometry}
\usepackage{hyperref}

\newcommand{\rk}{\operatorname{rank}}
\newcommand{\GL}{\operatorname{GL}}
\newcommand{\R}{\mathbb{R}}
\newcommand{\Mat}{\mathrm{Mat}}
\newcommand{\Hom}{\operatorname{Hom}}
\newcommand{\Rep}{\operatorname{Rep}}
\newcommand{\End}{\operatorname{End}}

\theoremstyle{definition}
\newtheorem{theorem}{Theorem}[section]
\newtheorem{proposition}[theorem]{Proposition}
\newtheorem{lemma}[theorem]{Lemma}
\newtheorem{corollary}[theorem]{Corollary}
\newtheorem{definition}[theorem]{Definition}
\newtheorem{remark}[theorem]{Remark}

\title{Task 1: Representations with Relations and $A_L^{\mathrm{wt}}$-Modules}
\date{}

\begin{document}
\maketitle

\section{Goal}

This file identifies the direct weight-space cyclic equations with the module
category of the algebra
\[
A_L^{\mathrm{wt}}=kC_{L+1}/I_L^{\mathrm{wt}}.
\]

The main deliverable is the correspondence:

\begin{itemize}
\item quiver representations satisfying the relations $p_\ell=0$;
\item modules over $A_L^{\mathrm{wt}}$;
\item after fixing the closing arrow $a_0=W_0$, points of the weight-space cyclic
  locus $\Gamma_{d,S}$.
\end{itemize}

\section{Representation spaces}

\begin{definition}[Representation of the cyclic quiver]
Fix a dimension vector
\[
d=(d_0,\dots,d_L),
\]
and vector spaces
\[
V_\ell\cong k^{d_\ell},
\qquad
0\le \ell\le L.
\]
The representation space of $Q=C_{L+1}$ is
\[
\Rep_d(Q)
:=
\prod_{\ell=1}^L \Hom(V_{\ell-1},V_\ell)\times \Hom(V_L,V_0).
\]
We write a point as
\[
\phi=(\phi_0,\phi_1,\dots,\phi_L).
\]
\end{definition}

\begin{definition}[Representations with weight-space relations]
Define the closed subvariety
\[
\Rep_d(Q,I_L^{\mathrm{wt}})
:=
\left\{
\phi\in \Rep_d(Q)
\;\middle|\;
p_\ell(\phi)=0 \text{ for every } 1\le \ell\le L
\right\}.
\]
\end{definition}

\begin{definition}[$A_L^{\mathrm{wt}}$-modules of dimension vector $d$]
Let
\[
V:=\bigoplus_{\ell=0}^L V_\ell.
\]
We denote by
\[
\Rep_d(A_L^{\mathrm{wt}})
\]
the set of $A_L^{\mathrm{wt}}$-module structures on $V$ for which the primitive
idempotent at vertex $\ell$ acts as the projection onto the summand $V_\ell$.
\end{definition}

\section{The standard correspondence}

\begin{proposition}[From representations with relations to $A_L^{\mathrm{wt}}$-modules]
There is a natural map
\[
F:\Rep_d(Q,I_L^{\mathrm{wt}})\to \Rep_d(A_L^{\mathrm{wt}})
\]
which sends a quiver representation $\phi$ to the unique
$A_L^{\mathrm{wt}}$-module whose arrow actions are the maps $\phi_\ell$.
\end{proposition}

\begin{proof}
Any quiver representation $\phi$ defines a $kQ$-module structure on
$V=\bigoplus_{\ell=0}^L V_\ell$ by letting each path act by the corresponding
composition of arrow maps.

If $\phi\in \Rep_d(Q,I_L^{\mathrm{wt}})$, then each generator $p_\ell$ acts by
zero. Since $I_L^{\mathrm{wt}}$ is generated by the $p_\ell$, the $kQ$-action
factors through the quotient algebra
\[
A_L^{\mathrm{wt}}=kQ/I_L^{\mathrm{wt}}.
\]
This gives the required $A_L^{\mathrm{wt}}$-module structure.
\end{proof}

\begin{proposition}[From $A_L^{\mathrm{wt}}$-modules to representations with relations]
There is a natural map
\[
G:\Rep_d(A_L^{\mathrm{wt}})\to \Rep_d(Q,I_L^{\mathrm{wt}})
\]
which sends an $A_L^{\mathrm{wt}}$-module to the quiver representation given by
the actions of the classes of the arrows $a_0,\dots,a_L$.
\end{proposition}

\begin{proof}
Given an $A_L^{\mathrm{wt}}$-module structure, precompose the action map
\[
kQ\twoheadrightarrow A_L^{\mathrm{wt}}\to \End_k(V)
\]
with the quotient morphism.
The classes of the arrows then act as linear maps
\[
\phi_\ell:V_{\ell-1}\to V_\ell,
\]
and every generator $p_\ell$ acts by zero in the quotient algebra. Hence the
resulting quiver representation lies in $\Rep_d(Q,I_L^{\mathrm{wt}})$.
\end{proof}

\begin{corollary}[The correspondence]
The maps $F$ and $G$ are mutually inverse bijections:
\[
\Rep_d(Q,I_L^{\mathrm{wt}})
\cong
\Rep_d(A_L^{\mathrm{wt}}).
\]
\end{corollary}

\section{Specializing the closing arrow}

\begin{definition}[Fixed-arrow representation slice]
Fix the data-dependent closing map
\[
W_0=\Sigma_{XY}P_Q\in \Hom(V_L,V_0).
\]
Define
\[
\Rep_{d,W_0}(Q)
:=
\{\phi\in \Rep_d(Q)\mid \phi_0=W_0\},
\]
and
\[
\Gamma_{d,S}^{\mathrm{quiv}}
:=
\Rep_{d,W_0}(Q)\cap \Rep_d(Q,I_L^{\mathrm{wt}}).
\]
\end{definition}

\begin{proposition}[Recovery of the direct weight-space cyclic locus]
Under the identification
\[
\phi_0=W_0,
\qquad
\phi_\ell=W_\ell
\quad
(1\le \ell\le L),
\]
the quiver slice $\Gamma_{d,S}^{\mathrm{quiv}}$ is naturally identified with the
direct weight-space cyclic locus
\[
\Gamma_{d,S}
=
\left\{
W=(W_1,\dots,W_L)
\;\middle|\;
W_{<\ell}W_0W_{>\ell}=0
\text{ for every } 1\le \ell\le L
\right\}.
\]
\end{proposition}

\begin{proof}
By Task 0, the relation
\[
p_\ell(\phi)=0
\]
is exactly the equation
\[
W_{<\ell}W_0W_{>\ell}=0.
\]
Therefore the condition defining $\Gamma_{d,S}^{\mathrm{quiv}}$ is identical to
the condition defining $\Gamma_{d,S}$ once the closing arrow is fixed to $W_0$.
\end{proof}

\begin{remark}
This is the module-theoretic reformulation needed for the second conjecture.
It says that the weight-space cyclic locus is not just an arbitrary algebraic
variety: after fixing $W_0$, it is the representation space of the algebra
$A_L^{\mathrm{wt}}$ inside a fixed-arrow slice.
\end{remark}

\section{Deliverables}

This task is complete once the file contains:

\begin{itemize}
\item the bijection
  $\Rep_d(Q,I_L^{\mathrm{wt}})\cong \Rep_d(A_L^{\mathrm{wt}})$;
\item the fixed-closing-arrow slice $\Gamma_{d,S}^{\mathrm{quiv}}$;
\item the identification
  $\Gamma_{d,S}^{\mathrm{quiv}}\cong \Gamma_{d,S}$.
\end{itemize}

\end{document}
